
\SGARefTarget{sga3:xvi:section:6}%
\section*{6. Applications: Extension of Certain Rigidity Properties
of Tori to Groups of Unipotent Rank Zero}
\addcontentsline{toc}{section}{6. Applications: Extension of Certain
Rigidity Properties of Tori to Groups of Unipotent Rank Zero}

\noindent\SGARefTarget{sga3:xvi:proposition:6:1}%
\textbf{Proposition 6.1.} --- \textit{Let \(S\) be a prescheme and
let \(G\) be an \(S\)-group prescheme smooth over \(S\), of finite
presentation, with connected affine fibers, and whose center is of
multiplicative type (for example, \(G\) of unipotent rank zero;
cf.\ \SGARefLink{sga3:xvi:theorem:4:1}{4.1} a)). Then every closed normal subgroup \(K\) of \(G\),
of finite presentation over \(S\) and quasi-finite over \(S\), is
finite over \(S\).}

\medskip
Indeed, for every geometric point \(\bar s\) above \(S\),
\((K_{\bar s})_{\mathrm{red}}\) is a finite étale subgroup of
\(G_{\bar s}\), normal and therefore central (since \(G_{\bar s}\)
is connected). Consequently
\[
  K'=Z\cap K,
\]
where \(Z\) denotes the center of \(G\), has the same underlying
space as \(K\), and it suffices to show that \(K'\) is finite over
\(S\). Now \(K'\) is a closed quasi-finite subgroup of the group
of multiplicative type \(Z\), and is therefore finite, as is
immediate (locally on \(S\), \(K'\) is contained in \({}_nZ\) for
a suitable integer \(n\), and \({}_nZ\) is finite over \(S\)).

\medskip
\noindent\SGARefTarget{sga3:xvi:corollary:6:2}%
\textbf{Corollary 6.2.} --- \textit{Let \(S\) and \(G\) be as
above, let \(K\) be an \(S\)-group subprescheme of \(G\), of
finite presentation over \(S\), normal and closed in \(G\), and
let \(s\) be a point of \(S\). If \(K_s\) is finite (respectively
is the unit group), there exists an open neighborhood \(U\) of
\(s\) such that \(K|_U\) is finite over \(S\) (respectively is the
unit group).}

\medskip
In view of the upper semicontinuity of the dimension of the fibers
of \(K\) (Exposé VI\textsubscript{B}, \S4), after shrinking \(S\)
we may already suppose that \(K\) is quasi-finite over \(S\), and
hence finite by \SGARefLink{sga3:xvi:proposition:6:1}{6.1}. If, moreover, \(K_s\) is the unit group, then
Nakayama's lemma easily implies that \(K\) is the unit group over
\(\operatorname{Spec}\mathcal O_s\), and hence over a neighborhood
of \(s\).

\medskip
\noindent\SGARefTarget{sga3:xvi:proposition:6:3}%
\textbf{Proposition 6.3.} --- \textit{Let \(S\) be a prescheme,
let \(G\) be an \(S\)-group prescheme smooth and of finite
presentation over \(S\), with affine, connected fibers, of
unipotent rank zero, let \(H\) be an \(S\)-group prescheme, let
\(s\) be a point of \(S\), and let}
\[
  u:G\longrightarrow H
\]
\textit{be an \(S\)-group morphism such that \(\Ker u_s\) is
central. Suppose, moreover, that \(H\) is of finite presentation
over \(S\), or that \(S\) is locally noetherian and \(G\) is
locally of finite type over \(S\). Then there exists an open
neighborhood \(U\) of \(s\) such that, if \(K=\Ker u\), \(K|_U\)
is central and of multiplicative type. Moreover,}
\[
  G'=(G|_U)/(K|_U)
\]
\textit{is representable, and the morphism}
\[
  u':G'\longrightarrow H|_U
\]
\textit{induced from \(u\) by passage to the quotient is an
immersion.}

\medskip
\noindent\textit{Proof.}
Let \(Z\) be the center of \(G\), and put \(K'=K\cap Z\).

\medskip
\noindent a) \(K'\) is a subgroup of multiplicative type over a
neighborhood \(U\) of \(s\). Indeed, after making an étale
extension over a neighborhood \(U\) of \(s\), which is legitimate,
we may suppose that \(G\) possesses a maximal torus \(T\)
(\SGARefLink{sga3:xvi:theorem:4:1}{4.1} c)). But then \(T\cap K\) is a group of multiplicative type
(\SGARefLink{sga3:xv:lemma:8:3}{Exposé XV, 8.3}) whose fiber at \(s\) is central by hypothesis;
therefore \(T\cap K\) is central over a neighborhood of \(s\)
(\SGARefLink{sga3:ix:theorem:5:6}{Exposé IX, 5.6} a)), and consequently coincides with \(K'\).

\medskip
\noindent b) Let us show that \(K'=K\) over a neighborhood \(U\)
of \(s\). By a), we may already suppose that \(K'\) is of
multiplicative type, and hence flat over \(S\). Since \(K'_s=K_s\),
the natural immersion
\[
  K'\longrightarrow K
\]
is then open over a neighborhood \(U\) of \(s\) (Exposé
VI\textsubscript{B}, \SGARefLink{cor:sga3-VIB-2.6}{2.6}). If \(t\in U\), the image of \(K_t\) in
the group
\[
  G'_t=G_t/Z_t
\]
is a finite étale group, normal in \(G'_t\) (since \(K_t\) is
normal in \(G_t\)), hence central, and therefore reduced to the
unit group (\SGARefLink{sga3:xvi:theorem:3:3}{3.3} b)). This says that \(K_t\) is contained in
\(Z_t\), and hence is equal to \(K'_t\), whence \(K=K'\) over
\(U\).

\medskip
\noindent c) The representability of \(G'\) now follows from \SGARefLink{sga3:xvi:corollary:2:3}{2.3};
the fact that \(u'\) is an immersion is contained in \SGARefLink{sga3:xvi:theorem:1:3}{1.3} a), in
view of \SGARefLink{sga3:xvi:theorem:4:1}{4.1} c) applied to \(G'\).

\medskip
\noindent\SGARefTarget{sga3:xvi:proposition:6:4}%
\textbf{Proposition 6.4.} --- \textit{Let \(S\) be a prescheme,
let \(G\) be an \(S\)-group prescheme smooth and of finite
presentation over \(S\), with connected fibers, and let \(s\) be
a point of \(S\) such that \(G_s\) is generated by its subtori
(\SGARefLink{sga3:xii:8.2}{Exposé XII, 8.2}) (for example, \(G_s\) affine of unipotent rank
zero). Let \(H\) be an \(S\)-group prescheme, and let}
\[
  u,v:G\longrightarrow H
\]
\textit{be two \(S\)-homomorphisms such that \(u_s=v_s\) and
\(u_s\) is central. Suppose, moreover, that \(H\) is of finite
presentation over \(S\), or that \(S\) is locally noetherian.
Then there exists a neighborhood \(U\) of \(s\) such that}
\[
  u|_U=v|_U.
\]

\medskip
We reduce as usual to the case where \(S\) is noetherian (to study
the condition ``\(G_s\) is generated by its subtori,'' one uses
Exposé VI\textsubscript{B}, \SGARefLink{prop:sga3-VIB-7.4}{7.4}). Since \(G\) has connected fibers,
to prove that \(u=v\) over a neighborhood \(U\) of \(s\), it
suffices to prove that \(u=v\) after reduction modulo every power
of the maximal ideal of the local ring \(\mathcal O_{S,s}\), which
reduces us to the case where \(S\) is local artinian.

But then the functor of maximal subtori of \(G\) is representable
by an \(S\)-scheme \(X\), smooth and of finite type over \(S\)
(\SGARefLink{sga3:xv:corollary:8:17}{Exposé XV, 8.17}). Let \(T\) be the maximal subtorus of \(G_X\)
universal for the functor \(X\). The hypothesis on \(G_s\) means
that the algebraic subgroup of \(G_s\) generated (Exposé
VI\textsubscript{B}, \S7) by the \(\kappa(s)\)-morphism
\[
  f:T_s\longrightarrow G_s
\]
(the composite of the immersion \(T_s\to G_{X_s}\) and the
canonical projection \(G_{X_s}\to G_s\)) is all of \(G_s\). But
\(X_s\) is geometrically connected (for if \(\overline T\) is a
maximal torus of \(G_{\bar s}\), and \(N\) its normalizer in
\(G_{\bar s}\), then \(X_{\bar s}\) is isomorphic to
\(G_{\bar s}/N\)), and the image of \(T_s\) under \(f\) contains
the unit section of \(G_s\). It follows from Exposé
VI\textsubscript{B}, \SGARefLink{prop:sga3-VIB-7.4}{7.4}, that there exist an integer \(N>0\) and
an \(S\)-morphism
\[
  f^N:T^N\longrightarrow G,
  \qquad
  T^N=T\times_S\cdots\times_S T
  \quad\text{(\(N\) factors),}
\]
where \(f^N\) depends only on \(f\) and the multiplication law of
\(G\), such that \((f^N)_s\) is surjective.

Consider the base change \(X\to S\) and the restrictions of
\(u_X\) and \(v_X\) to the subtorus \(T\) of \(G_X\). By
hypothesis,
\[
  u_{X_s}=v_{X_s},
\]
and
\[
  u_{X_s}:T_{X_s}\longrightarrow H_{X_s}
\]
is a central homomorphism. It then follows from \SGARefLink{sga3:ix:theorem:5:1}{Exposé IX, 5.1},
that
\[
  u_X|_T=v_X|_T.
\]
Making the definition of \(f^N\) explicit, and taking account of
the fact that \(u\) and \(v\) are homomorphisms, one immediately
deduces that
\[
  uf^N=vf^N.
\]
But the \(\kappa(s)\)-morphism \(f_s^N\) is surjective and \(G_s\)
is smooth, hence reduced; consequently \(f_s^N\) is generically
flat. Since \(T\) is smooth over \(S\), hence flat, there exists a
nonempty open subset \(V\) of \(T^N\) such that
\[
  f^N|_V
\]
is flat (EGA IV, 11.3.10).\footnote{The French source prints \(T\)
here; \(T^N\) is the domain of \(f^N\).} The image of \(V\) is an open subset
\(W\) of \(G\) (EGA IV, 11.3.1), and
\[
  f^N:V\longrightarrow W
\]
is faithfully flat and of finite presentation, hence covering for
the fpqc topology. The equality \(uf^N=vf^N\) then implies
\[
  u|_W=v|_W.
\]
Since \(W\) is schematically dense in \(G\) (EGA IV, 11.10.10),
and \(H\) is separated over \(S\) (Exposé VI\textsubscript{A},
\SGARefLink{sga3:VIA:0.3}{0.3}), it follows immediately that \(u=v\).
